\documentclass[10pt, letterpaper]{article}

% Packages:
\usepackage[
    ignoreheadfoot,
    top=0.4 cm,
    bottom=0.4 cm,
    left=0.8 cm,
    right=0.8 cm,
    footskip=0.2 cm,
]{geometry}
\usepackage{titlesec}
\usepackage{tabularx}
\usepackage{array}
\usepackage[dvipsnames]{xcolor}
\definecolor{primaryColor}{RGB}{0, 0, 0}
\usepackage{enumitem}
\usepackage{amsmath}
\usepackage[
    pdftitle={Tanmay Bhuskute's CV},
    pdfauthor={Tanmay Bhuskute},
    pdfcreator={LaTeX with RenderCV},
    colorlinks=true,
    urlcolor=primaryColor
]{hyperref}
\usepackage[pscoord]{eso-pic}
\usepackage{calc}
\usepackage{bookmark}
\usepackage{lastpage}
\usepackage{changepage}
\usepackage{paracol}
\usepackage{ifthen}
\usepackage{needspace}
\usepackage{iftex}

\ifPDFTeX
    \input{glyphtounicode}
    \pdfgentounicode=1
    \usepackage[T1]{fontenc}
    \usepackage[utf8]{inputenc}
    \usepackage{lmodern}
\fi

\usepackage{XCharter}

% Some settings:
\raggedright
\AtBeginEnvironment{adjustwidth}{\partopsep0pt}
\pagestyle{empty}
\setcounter{secnumdepth}{0}
\setlength{\parindent}{0pt}
\setlength{\topskip}{0pt}
\setlength{\columnsep}{0.15cm}
\pagenumbering{gobble}

% ── Section title ──
\titleformat{\section}{\needspace{2\baselineskip}\bfseries\large}{}{0pt}{}[\vspace{1pt}\titlerule]
\titlespacing{\section}{-1pt}{0.15 cm}{0.08 cm}

\renewcommand\labelitemi{$\vcenter{\hbox{\small$\bullet$}}$}

% ── Bullet lists ──
\newenvironment{highlights}{
    \begin{itemize}[
        topsep=0.10 cm,
        parsep=0.06 cm,
        partopsep=0pt,
        itemsep=0.06 cm,
        leftmargin=12pt
    ]
}{
    \end{itemize}
}

\newenvironment{highlightsforbulletentries}{
    \begin{itemize}[
        topsep=0.10 cm,
        parsep=0.06 cm,
        partopsep=0pt,
        itemsep=0.04 cm,
        leftmargin=12pt
    ]
}{
    \end{itemize}
}

\newenvironment{onecolentry}{
    \begin{adjustwidth}{0 cm + 0.00001 cm}{0 cm + 0.00001 cm}
}{
    \end{adjustwidth}
}

\newenvironment{twocolentry}[2][]{
    \onecolentry
    \def\secondColumn{#2}
    \setcolumnwidth{\fill, 4.5 cm}
    \begin{paracol}{2}
}{
    \switchcolumn \raggedleft \secondColumn
    \end{paracol}
    \endonecolentry
}

\newenvironment{threecolentry}[3][]{
    \onecolentry
    \def\thirdColumn{#3}
    \setcolumnwidth{, \fill, 4.5 cm}
    \begin{paracol}{3}
    {\raggedright #2} \switchcolumn
}{
    \switchcolumn \raggedleft \thirdColumn
    \end{paracol}
    \endonecolentry
}

\newenvironment{header}{
    \setlength{\topsep}{0pt}\par\kern\topsep\centering\linespread{0.6}
}{
    \par\kern\topsep
}

\let\hrefWithoutArrow\href

\begin{document}
    \newcommand{\AND}{\unskip
        \cleaders\copy\ANDbox\hskip\wd\ANDbox
        \ignorespaces
    }
    \newsavebox\ANDbox
    \sbox\ANDbox{$|$}

    \begin{header}
        \fontsize{16 pt}{24 pt}\selectfont Tanmay Bhuskute

        \vspace{3 pt}

        \normalsize
        \mbox{ +1 (732) 858-3027 }%
        \kern 1.0 pt%
        \AND%
        \kern 1.0 pt%
        \mbox{ tanmay.v.bhuskute@gmail.com }%
        \kern 1.0 pt%
        \AND%
        \kern 1.0 pt%
        \mbox{ github.com/Tanmayb05 }%
        \kern 1.0 pt%
        \AND%
        \kern 1.0 pt%
        \mbox{ linkedin.com/in/tanmay-bhuskute }%
    \end{header}

    \vspace{2 pt}

%═══════════════════════════════════════
\section{SUMMARY}
%═══════════════════════════════════════
\begin{onecolentry}
Software Engineer and M.S. Computer Science candidate (Dec 2026) with production Python engineering experience building backend services, REST APIs, and AI agent workflows using Python, Boto3, AWS, and Kubernetes. Current research in LLM-powered agentic pipelines with prompt engineering and tool-calling.
\end{onecolentry}

%═══════════════════════════════════════
\section{EDUCATION}
%═══════════════════════════════════════
\begin{twocolentry}{
    Aug 2024 - Dec 2026
}
    \textbf{Master of Science in Computer Science} \\
    University of Cincinnati
\end{twocolentry}
\begin{onecolentry}
Relevant Coursework: Advanced Algorithms, Deep Learning, Intelligent Data Analysis, Database Theory.
\end{onecolentry}
\vspace{0.15 cm}

\begin{twocolentry}{
    Aug 2018 - Jul 2022
}
    \textbf{Bachelor of Engineering in Computer Engineering} \\
    AISSMS College of Engineering
\end{twocolentry}
\begin{onecolentry}
Relevant Coursework: Data Mining, Data Analytics, Machine Learning, Object-Oriented Programming.
\end{onecolentry}
%═══════════════════════════════════════
\section{TECHNICAL SKILLS}
%═══════════════════════════════════════
\begin{onecolentry}
\textbf{Programming:} Python, Java, JavaScript, TypeScript, C/C++, SQL, Go.
\end{onecolentry}
\vspace{0.06 cm}
\begin{onecolentry}
\textbf{Backend \& APIs:} REST APIs, Django, Flask, FastAPI, backend service architecture, data structures, React, Node.js.
\end{onecolentry}
\vspace{0.06 cm}
\begin{onecolentry}
\textbf{Testing \& CI/CD:} Pytest, JUnit, Git, GitLab CI/CD, Jenkins, unit/integration testing, Agile/Scrum.
\end{onecolentry}
\vspace{0.06 cm}
\begin{onecolentry}
\textbf{Developer Tools:} AWS, Docker, Kubernetes, ArgoCD, Helm, Terraform, Linux, Bash, SonarQube, Azure, Google Cloud Platform.
\end{onecolentry}
\vspace{0.06 cm}
\begin{onecolentry}
\textbf{AI \& Data:} Gemini, Vertex AI, prompt engineering, agentic workflows, LangChain, PyTorch, TensorFlow, pandas, NumPy.
\end{onecolentry}

%═══════════════════════════════════════
\section{WORK EXPERIENCE}
%═══════════════════════════════════════
\begin{twocolentry}{Feb 2026 -- Present}
\textbf{Graduate Researcher} ---- University of Cincinnati
\end{twocolentry}
\begin{highlights}
  \item Built agentic LLM pipeline in Python with prompt engineering and tool-calling for Android GUI bug reproduction, debugging LLM outputs against \textbf{ViBR} baseline (GPT-4o + GroundingDINO + CLIP) across 53 runs on 20+ apps.
  \item Demonstrated \textbf{25.7x latency reduction} (39.4s vs 1012.1s avg per run) while preserving task-completion fidelity through iterative debugging of agent decision loops.
\end{highlights}

\vspace{0.15 cm}

\begin{twocolentry}{Jul 2022 -- Jun 2024}
\textbf{Software Development Engineer} ---- Siemens Industry Software
\end{twocolentry}
\begin{highlights}
  \item Wrote clean, maintainable Python (Boto3) backend automation for Teamcenter X infrastructure workflows, replacing legacy Ruby scripts and reducing deployment time by \textbf{40\%} across AWS environments.
  \item Debugged and resolved \textbf{44 production environment} issues in 4 months across Linux, EC2, databases, and backend services, using Git workflows and CI/CD pipelines for rapid, reliable fixes.
  \item Designed Ansible-based Blue-Green deployment framework, reducing production upgrade downtime by \textbf{50\%}.
  \item Migrated \textbf{40+ production environments} to Kubernetes using ArgoCD, Helm, and GitLab CI/CD, participating in sprint planning and retrospectives within an Agile Release Train.
  \item Served as Quality Coach and SonarQube Specialist, implementing pytest-based automated testing that raised code coverage from \textbf{35\% to 78\%}; led code reviews mentoring \textbf{5 software engineers}.
\end{highlights}

\vspace{0.15 cm}

\begin{twocolentry}{Mar 2022 -- Jun 2022}
\textbf{Software Developer Intern} ---- Siemens Industry Software
\end{twocolentry}
\begin{highlights}
  \item Automated provisioning of AWS resources with \textbf{Terraform}, reducing deployment time by \textbf{30\%}.
  \item Developed \textbf{Go + Terratest} test framework, catching \textbf{95\%} of misconfigurations pre-production.
\end{highlights}

\vspace{0.15 cm}

\begin{twocolentry}{Sep 2021 -- Dec 2021}
\textbf{Software Developer Intern} ---- Government of Maharashtra, Water Resource Department
\end{twocolentry}
\begin{highlights}
  \item Built and deployed a full-stack \textbf{React and MySQL} workflow application that digitized departmental record management, automated manual processes, and improved operational efficiency by \textbf{40\%}.
\end{highlights}

%═══════════════════════════════════════
\section{PROJECTS AND RESEARCH}
%═══════════════════════════════════════
\begin{twocolentry}{May 2026}
\textbf{PolicyPilot -- Medical Benefit Policy Intelligence Portal} [Python, Node, Gemini, MCP, REST APIs]
\end{twocolentry}
\begin{highlights}
  \item Designed \textbf{17-tool MCP server} in Python exposing structured LLM tool-calling and agentic workflows for PA extraction, patient-readiness evaluation, and cross-payer comparison, consumable by external AI agents.
  \item Architected evidence-grounded \textbf{RAG pipeline} over Google Gemini 2.5-flash with REST API backend, forcing citation-backed answers to eliminate hallucination in high-stakes healthcare decisions.
\end{highlights}

\vspace{0.12 cm}

\begin{twocolentry}{Jul 2026}
\textbf{NYC Mobility Forecasting} [Python, PySpark, scikit-learn, FastAPI]
\end{twocolentry}
\begin{highlights}
  \item Built Python ML system forecasting NYC taxi/rideshare demand across \textbf{263 zones} (\textbf{2.3M-row} dataset) served via a FastAPI REST backend; improved MAPE \textbf{18.77\% to 16.52\%} and cut RMSE \textbf{43\%} over baseline.
\end{highlights}

%═══════════════════════════════════════
\section{EXTRACURRICULAR ACTIVITIES}
%═══════════════════════════════════════
\begin{onecolentry}
\begin{highlights}
  \item \textbf{Finalist, RevolutionUC 2025 Hackathon} (300 participants, 31 schools) \,|\, \textbf{SAFe Practitioner 5.1} (Mar 2023).
\end{highlights}
\end{onecolentry}

\end{document}
